\chapter{Strategia di verifica}
\label{cap:td-strategia}

\section{Livelli di verifica}

La verifica è organizzata su tre livelli, con obiettivi distinti e costi
crescenti.

\begin{longtable}{@{}L{2.8cm} L{4.4cm} L{6.8cm}@{}}
\toprule
\textbf{Livello} & \textbf{Che cosa verifica} & \textbf{Come} \\
\midrule
\endfirsthead
\toprule
\textbf{Livello} & \textbf{Che cosa verifica} & \textbf{Come} \\
\midrule
\endhead
\bottomrule
\endlastfoot

Unità &
Le regole di dominio dei servizi e le funzioni di infrastruttura prese
singolarmente. &
Il servizio è costruito con DAO simulati passati al costruttore. Nessuna
rete, nessun database. \\

Integrazione &
La composizione fra rotte, middleware e serializzazione: controllo degli
accessi, validazione, forma delle risposte, traduzione degli errori. &
L'applicazione Express è costruita con servizi simulati e interrogata via
HTTP con \id{supertest}, senza occupare una porta. \\

Sistema &
Il funzionamento end-to-end su un'istanza reale, database compreso. &
Prove manuali documentate in \cref{sec:td-sistema}, eseguite su una base di
dati creata dagli script del progetto. \\

\end{longtable}

Il livello di unità e quello di integrazione sono automatici ed eseguiti a
ogni modifica del codice. Il livello di sistema è manuale: è l'unico che
esercita le interrogazioni SQL reali, ed è quindi la contromisura al
compromesso dichiarato nel SDD sulla sostituzione del livello di persistenza.

\section{Approccio}

L'approccio è \term{white-box}: i casi di test sono derivati conoscendo la
struttura interna del codice, e la loro adeguatezza è misurata sulla
copertura raggiunta. La scelta è motivata dal fatto che gli oracoli più
importanti del sistema — la soglia di superamento di un quiz, l'esclusione
delle soluzioni dalla serializzazione, la propagazione della connessione
transazionale — non sono osservabili dall'esterno senza conoscere il
progetto.

I casi di test non sono però scritti guardando l'implementazione riga per
riga: sono derivati dalle specifiche con il metodo della category partition
(\cref{cap:td-derivazione}), e la conoscenza della struttura interna serve a
verificare che nessun percorso sia rimasto scoperto. Un test scritto per
inseguire una riga di codice verifica che il codice faccia ciò che fa, non
ciò che deve fare.

\section{Criteri di ingresso e di uscita}

\begin{longtable}{@{}L{4.0cm} L{10.4cm}@{}}
\toprule
\textbf{Criterio} & \textbf{Condizione} \\
\midrule
\endfirsthead
\toprule
\textbf{Criterio} & \textbf{Condizione} \\
\midrule
\endhead
\bottomrule
\endlastfoot

Ingresso &
Il codice compila senza errori di tipo e supera l'analisi statica. Una
funzionalità non entra in verifica finché la sua specifica non è formulata
come requisito con criterio di accettazione. \\

Uscita &
Tutti i casi di test superati; copertura non inferiore all'85\% delle
istruzioni e all'80\% delle diramazioni; nessun difetto aperto di gravità
alta o critica. \\

\end{longtable}

Le soglie di copertura non sono un valore riportato a posteriori: sono
configurate nella suite, che fallisce se non vengono raggiunte. La differenza
è sostanziale — una soglia dichiarata in un documento è un auspicio, una
soglia imposta dalla suite è un vincolo.

\section{Sospensione e ripresa}

L'esecuzione della verifica è sospesa quando il sistema non soddisfa più i
criteri d'ingresso: un errore di compilazione, un fallimento dell'analisi
statica o un difetto che impedisce l'esecuzione della suite. In questi casi
proseguire non produrrebbe informazione, perché i fallimenti successivi
sarebbero conseguenza del primo.

La verifica riprende quando la causa è rimossa e la suite viene rieseguita per
intero. Non si riprende dall'ultimo test fallito: un difetto corretto può
averne introdotto un altro altrove, e solo un'esecuzione completa lo
rivelerebbe.

\section{Ambiente e strumenti}

\begin{longtable}{@{}L{3.4cm} L{11.0cm}@{}}
\toprule
\textbf{Strumento} & \textbf{Ruolo} \\
\midrule
\endfirsthead
\toprule
\textbf{Strumento} & \textbf{Ruolo} \\
\midrule
\endhead
\bottomrule
\endlastfoot

Jest &
Esecutore dei test, gestione dei \term{mock} e misurazione della copertura. \\

ts-jest &
Compilazione TypeScript durante l'esecuzione: i test sono scritti nello
stesso linguaggio del codice e ne condividono i tipi. \\

supertest &
Invio di richieste HTTP all'applicazione Express in memoria, senza aprire
una porta di rete. \\

TypeScript &
Controllo dei tipi in modalità stretta, primo livello di verifica: molte
classi di errore non arrivano mai a essere eseguite. \\

ESLint e Prettier &
Analisi statica e uniformità formale del codice. \\

GitHub Actions &
Esecuzione automatica di analisi statica, controllo dei tipi, test e
compilazione a ogni integrazione. \\

\end{longtable}

\subsection{Isolamento dal database}

Il modulo che espone il pool di connessioni è sostituito globalmente da un
mock all'avvio della suite. Nessun test apre una connessione: l'intera
verifica automatica gira su qualunque macchina senza MySQL installato e senza
variabili d'ambiente reali.

La scelta ha due conseguenze da tenere presenti. La prima è positiva: la
suite completa impiega pochi secondi, quindi viene eseguita spesso. La
seconda è il limite già dichiarato: le interrogazioni SQL non sono verificate
automaticamente (\cref{sec:td-limiti}).

\section{Funzionalità da verificare}

\begin{longtable}{@{}L{4.4cm} L{10.0cm}@{}}
\toprule
\textbf{Area} & \textbf{Oggetto della verifica} \\
\midrule
\endfirsthead
\toprule
\textbf{Area} & \textbf{Oggetto della verifica} \\
\midrule
\endhead
\bottomrule
\endlastfoot

Autenticazione &
Politica delle password, hashing, riconoscimento delle credenziali,
indistinguibilità fra email inesistente e password errata, apertura e
chiusura della sessione. \\

Gestione account &
Aggiornamento parziale del profilo, unicità dell'email, cambio password con
verifica di quella attuale, cancellazione. \\

Gestione tutorial &
Vincoli di dominio su titolo, testo, categoria e immagine; sanificazione del
contenuto; filtro e ricerca; ciclo di vita del tutorial. \\

Gestione quiz &
Vincoli sulle domande; correzione e soglia di superamento; associazione delle
risposte per identificativo; conservazione di un esito positivo; esclusione
delle soluzioni dalla rappresentazione. \\

Gestione feedback &
Vincoli su valutazione e commento; unicità per coppia utente-tutorial;
distinzione fra proprietà del dato e potere di moderazione. \\

Gestione obiettivi &
Assegnazione automatica alle soglie; idempotenza; abbinamento fra
conseguimenti e obiettivi. \\

Controllo degli accessi &
Matrice completa: per ogni rotta protetta, comportamento verso anonimo,
utente e amministratore. \\

Infrastruttura &
Gerarchia degli errori, traduzione in risposte HTTP, gestione dei nomi di
file, composizione dell'applicazione. \\

\end{longtable}

\section{Funzionalità non verificate automaticamente}

\begin{longtable}{@{}L{4.4cm} L{10.0cm}@{}}
\toprule
\textbf{Esclusione} & \textbf{Motivazione} \\
\midrule
\endfirsthead
\toprule
\textbf{Esclusione} & \textbf{Motivazione} \\
\midrule
\endhead
\bottomrule
\endlastfoot

Interrogazioni SQL e vincoli del database &
Il livello di persistenza è sostituito da mock. Coperti dalle prove di
sistema (\cref{sec:td-sistema}). \\

Trigger di ricalcolo della valutazione media &
Vivono nel database. Coperti dalle prove di sistema. \\

Interfaccia utente &
Non esistono test automatici sul client. Il comportamento è verificato
manualmente; il controllo dei tipi e l'analisi statica coprono la parte
strutturale. \\

Elaborazione delle immagini &
Il ridimensionamento è delegato a una libreria esterna e non è
riverificato. La logica applicativa che la circonda — generazione del nome,
validazione del percorso — è invece coperta. \\

Widget conversazionale &
Componente esterno, fuori dal perimetro del sistema. \\

\end{longtable}
